\documentclass[11pt]{article}

\usepackage[margin=1in]{geometry}
\usepackage{amsmath,amssymb,amsthm,mathtools}
\usepackage{hyperref}
\usepackage{enumitem}
\usepackage{xcolor}

\newtheorem{theorem}{Theorem}
\newtheorem{lemma}{Lemma}
\newtheorem{proposition}{Proposition}
\newtheorem{conjecture}{Bridge Target}
\newtheorem{definition}{Definition}
\newtheorem{remark}{Remark}

\newcommand{\C}{\mathbb C}
\newcommand{\R}{\mathbb R}
\newcommand{\zetaR}{\zeta}
\newcommand{\NTZ}{\mathcal Z_{\mathrm{nt}}}
\newcommand{\OLZ}{\mathcal Z_{\mathrm{off}}}
\newcommand{\CL}{\{s\in\C:\Re(s)=1/2\}}
\newcommand{\gpdef}{\mathcal D_G}
\newcommand{\KL}{K_L}
\newcommand{\KR}{K_R}

\title{A Cosh--Weil Detector Route to the Riemann Hypothesis\\
\large Draft proof architecture from the Lean formalization}
\author{Sam Lavery}
\date{\today}

\begin{document}
\maketitle

\begin{abstract}
This note describes the proof architecture currently represented in the Lean repository
\texttt{RequestProject}.  The argument separates into three layers.  First, a purely
geometric cosh/Gaussian detector proves that a real parameter \(\beta\ne 1/2\) carries a
strictly positive pair-defect.  Second, the zeta-zero layer defines nontrivial and off-line
zeros and translates vanishing of that pair-defect at every nontrivial zero into the
statement that every nontrivial zero lies on the critical line.  Third, a Weil explicit
formula layer is intended to prove the analytic vanishing/balance target for the chosen
pair-cosh Gaussian test.  The present document treats the first and third layers with
different status labels: the cosh detector facts are formal geometric consequences, while
the Weil vanishing target is the load-bearing analytic bridge unless the corresponding
Lean theorem has been independently discharged.
\end{abstract}

\section{Lean-facing statement of the goal}

Let
\[
  \NTZ := \{\rho\in\C : 0<\Re(\rho)<1 \text{ and } \zetaR(\rho)=0\}.
\]
The off-line zeros are
\[
  \OLZ := \{\rho\in\NTZ : \Re(\rho)\ne 1/2\}.
\]
The formal endgame is to prove
\[
  \forall \rho\in\NTZ,\quad \Re(\rho)=1/2,
\]
and then invoke the bridge to Mathlib's literal \texttt{RiemannHypothesis} predicate.

\begin{theorem}[RH bridge]
If every \(\rho\in\NTZ\) satisfies \(\Re(\rho)=1/2\), then Mathlib's
\texttt{RiemannHypothesis} proposition follows.
\end{theorem}

\begin{remark}
In Lean this is represented by the bridge theorem
\texttt{RHBridge.no\_offline\_zeros\_implies\_rh}.  Its role is not to prove the
analytic zero placement by itself; it upgrades the internal zero-set statement to
Mathlib's global formulation, including the standard exclusions of trivial zeros and the
pole at \(s=1\).
\end{remark}

\section{The detector-pair geometry}

The basic two-kernel construction uses the two reflected anchors
\[
  a=\frac{\pi}{6},\qquad 1-a=1-\frac{\pi}{6}.
\]
For real \(\beta,t\), define
\[
  \KL(\beta,t)=\cosh\bigl((\beta-a)t\bigr),\qquad
  \KR(\beta,t)=\cosh\bigl((\beta-(1-a))t\bigr).
\]
The reflection \(\beta\mapsto 1-\beta\) exchanges the two detectors:
\[
  \KL(1-\beta,t)=\KR(\beta,t),\qquad
  \KR(1-\beta,t)=\KL(\beta,t).
\]

\begin{lemma}[Detector agreement classifier]
For \(t\ne 0\),
\[
  \KL(\beta,t)=\KR(\beta,t) \quad\Longleftrightarrow\quad \beta=\frac12.
\]
\end{lemma}

\begin{proof}[Proof idea]
The equality \(\cosh x=\cosh y\) forces \(|x|=|y|\).  With
\(x=(\beta-a)t\) and \(y=(\beta-(1-a))t\), the branch \(x=-y\) gives
\((2\beta-1)t=0\), hence \(\beta=1/2\).  The branch \(x=y\) would force
\((1-2a)t=0\), i.e. \((1-\pi/3)t=0\), impossible for \(t\ne0\) because
\(\pi\ne3\).
\end{proof}

\section{Gaussian pair defect}

Let \(\psi(t)=e^{-t^2}\).  The pointwise squared detector difference is
\[
  \Delta(\beta,t)^2 := \bigl(\KL(\beta,t)-\KR(\beta,t)\bigr)^2.
\]
The key factorization is
\[
  \Delta(\beta,t)^2
  =4\sinh^2\!\left(\left(\frac12-\frac\pi6\right)t\right)
     \sinh^2\!\left(\left(\beta-\frac12\right)t\right).
\]
Define the Gaussian pair defect
\[
  \gpdef(\beta)
  :=\int_0^\infty \Delta(\beta,t)^2\psi(t)^2\,dt.
\]

\begin{lemma}[Vanishing on the line]
\[
  \gpdef(1/2)=0.
\]
\end{lemma}

\begin{proof}[Proof idea]
The factor \(\sinh((\beta-1/2)t)\) vanishes identically when \(\beta=1/2\).
\end{proof}

\begin{lemma}[Nonnegativity]
For all \(\beta\in\R\),
\[
  \gpdef(\beta)\ge0.
\]
\end{lemma}

\begin{proof}[Proof idea]
The integrand is a square multiplied by the nonnegative Gaussian weight.
\end{proof}

\begin{lemma}[Strict off-line positivity]
If \(\beta\ne1/2\), then
\[
  \gpdef(\beta)>0.
\]
\end{lemma}

\begin{proof}[Proof idea]
For every \(t>0\), both constants
\(1/2-\pi/6\) and \(\beta-1/2\) are nonzero.  Hence both sinh factors are
nonzero on \((0,\infty)\), and the Gaussian weight is strictly positive.  The
integrand is therefore positive on a set of positive measure, giving a strictly
positive integral.
\end{proof}

\begin{proposition}[Defect zero forces the critical line]
For every real \(\beta\),
\[
  \gpdef(\beta)=0 \quad\Longrightarrow\quad \beta=1/2.
\]
\end{proposition}

\begin{proof}
The contrapositive is exactly strict off-line positivity.
\end{proof}

\section{The zeta-zero bridge target}

The detector geometry becomes an RH argument once one proves that the Gaussian pair
defect vanishes at the real part of every nontrivial zeta zero.

\begin{conjecture}[Pair Gaussian bridge / Weil vanishing target]
For every nontrivial zeta zero \(\rho\in\NTZ\),
\[
  \gpdef(\Re(\rho))=0.
\]
\end{conjecture}

\begin{theorem}[Conditional RH from pair Gaussian bridge]
Assume the pair Gaussian bridge.  Then every nontrivial zero lies on the critical line:
\[
  \forall \rho\in\NTZ,\quad \Re(\rho)=1/2.
\]
Consequently Mathlib's \texttt{RiemannHypothesis} follows via the RH bridge.
\end{theorem}

\begin{proof}
Let \(\rho\in\NTZ\).  The bridge target gives \(\gpdef(\Re\rho)=0\).  The defect
classifier gives \(\Re\rho=1/2\).  Since \(\rho\) was arbitrary, all nontrivial
zeros are on the critical line.  The Lean bridge then upgrades this internal statement
to Mathlib's \texttt{RiemannHypothesis}.
\end{proof}

\section{Even and odd channel interpretation}

The formalization also packages the pair defect as an even-channel observable:
\[
  \mathrm{Even}(\beta):=\gpdef(\beta).
\]
It is nonnegative, vanishes at \(1/2\), and is strictly positive off the line.
The odd channel is a signed phase-sensitive companion of the form
\[
  \mathrm{Odd}(\beta,\gamma)
  =\int_{\R}2\sinh((\beta-1/2)t)\sin(\gamma t)\psi(t)^2\,dt.
\]
It has the closed form
\[
  \mathrm{Odd}(\beta,\gamma)
  =\sqrt{2\pi}\exp\left(\frac{(\beta-1/2)^2-\gamma^2}{8}\right)
    \sin\left(\frac{(\beta-1/2)\gamma}{4}\right).
\]
The odd channel is anti-invariant under \(\beta\mapsto1-\beta\), while the even channel
is invariant.  This supports the Klein-orbit language: the sine/odd contribution cancels
under the functional-equation orbit, while the even/cosh excess does not cancel unless it
is zero.

\section{Weil explicit formula layer}

The intended analytic bridge is a Weil explicit formula for the pair-cosh Gaussian test.
The repository's final assembly scaffold decomposes this into rectangle contour pieces:

a. rectangle Cauchy--Goursat decomposition;

b. residues at nontrivial zeros with multiplicity;

c. horizontal edge decay on good heights;

d. right-edge conversion to an Euler-product/prime sum;

e. left-edge transfer by the functional equation into archimedean and reflected-prime
terms;

f. trivial-zero residue handling;

g. final assembly and extraction of per-zero positivity/vanishing.

The proof-critical analytic output should be recorded in a single theorem shape:
\[
  \forall \rho\in\NTZ,\quad \gpdef(\Re\rho)=0.
\]
Everything before this point is detector geometry; everything proving this statement is
Weil/explicit-formula analysis.

\section{Recommended theorem map}

A clean paper/Lean correspondence should use the following dependency order.

\begin{enumerate}[label=\textbf{T\arabic*.},leftmargin=3em]
  \item Define \(\NTZ\), \(\OLZ\), and the internal critical-line statement.
  \item Prove the detector swap \(\KL(1-\beta,t)=\KR(\beta,t)\).
  \item Prove the detector agreement classifier \(\KL=\KR\iff\beta=1/2\) for \(t\ne0\).
  \item Prove the squared-difference sinh factorization.
  \item Define \(\gpdef\) and prove \(\gpdef(1/2)=0\), \(\gpdef\ge0\), and
        \(\beta\ne1/2\Rightarrow\gpdef(\beta)>0\).
  \item State and prove the Weil pair Gaussian bridge:
        \(\forall\rho\in\NTZ,\gpdef(\Re\rho)=0\).
  \item Combine T5 and T6 to get \(\forall\rho\in\NTZ,\Re\rho=1/2\).
  \item Invoke the Mathlib bridge to obtain \texttt{RiemannHypothesis}.
\end{enumerate}

\section{Current audit notes}

The detector side is structurally strong: it is a direct real-variable argument using
cosh symmetry, sinh factorization, positivity of squares, and Gaussian positivity.

The analytic side should be audited for exactly one point: whether the repository already
contains a theorem with no hypotheses and no hidden assumptions proving the pair Gaussian
bridge.  If the available theorem is instead a named 	exttt{Prop} target or a theorem
conditional on final-assembly hypotheses, the proof is not yet closed; it is a precise
conditional RH reduction.

The remaining uniqueness/totality obligations do not appear to require new breakthrough
mathematics.  They are mostly Lean formalization grind: pushing through Fubini exchange,
absolute convergence estimates, analytic continuation of the cosh-transform family,
Riemann--Lebesgue decay on imaginary arguments, Fourier-cosine injectivity, and the
countable 	exttt{tsum} layer-peeling/moment-uniqueness argument.  These are serious
engineering lemmas because Mathlib does not package them in exactly the needed form, but
the intended mathematics is standard analysis plus careful summability bookkeeping rather
than a new RH-specific idea.

\section{Minimal final theorem shape}

The final Lean theorem should have the shape
\[
\texttt{theorem rh\_from\_pair\_weil : RiemannHypothesis := by \dots}
\]
with no bridge hypotheses in the theorem signature.  Internally, it should instantiate:
\[
  \texttt{PairGaussianBridge}
\]
from the completed Weil explicit formula, then apply the detector classifier and finally
\texttt{RHBridge.no\_offline\_zeros\_implies\_rh}.

\end{document}
